\chapter{Einleitung}

\section{Das Problem der modernen Physik}

Die moderne Physik befindet sich in einer paradoxen Lage. Einerseits verfügt sie mit der Quantenmechanik und der Allgemeinen Relativitätstheorie über zwei der erfolgreichsten Theorien der Wissenschaftsgeschichte. Beide liefern präzise Vorhersagen, beide sind experimentell bestätigt, und doch stehen sie in einem fundamentalen Widerspruch zueinander.

Die Quantenmechanik ist nicht-lokal, probabilistisch und basiert auf einer Wellenfunktion im Konfigurationsraum. Die Allgemeine Relativitätstheorie ist lokal, deterministisch und beschreibt Gravitation als Krümmung der Raumzeit. Beide Theorien funktionieren – aber sie können nicht gleichzeitig fundamental sein.

Die vorgeschlagenen Lösungen – Stringtheorie, Schleifenquantengravitation, dunkle Materie, dunkle Energie – haben das Problem nicht gelöst, sondern die Theorielandschaft weiter verkompliziert. Es fehlt ein einheitliches, konsistentes Fundament.

\section{Die Lösung: IWT und Q}

Dieses Dokument stellt die Informations-Weber-Theorie (IWT) vor.

Sie hat nur eine fundamentale Größe: die Information \( I \).

Aus der Information emergiert alles andere: Raum, Zeit, Materie, Energie, Gravitation, Quantenmechanik, Elektrodynamik.

Die organisierende Kraft ist \( Q \): die Quelle der Selbstorganisation.

\section{Zwei Quellen}

Dieses Dokument ruht auf zwei Säulen:

\begin{enumerate}
\item \textbf{Die Theorie:} Axiome, Herleitungen, mathematische Struktur.
\item \textbf{Die Simulation:} Numerische Bestätigung, Ergebnisse, Visualisierung.
\end{enumerate}

Die Theorie sagt, was sein muss. Die Simulation zeigt, dass es ist.

\section{Die zentrale Aussage}

\( Q \) ist die Quelle.

Alles emergiert aus \( Q \).

\begin{itemize}
\item Die DBT emergiert aus \( Q \) (Bohm-Potential).
\item Die WG emergiert aus \( Q \) (Gravitation).
\item Die WED emergiert aus \( Q \) (Elektrodynamik).
\end{itemize}

\( Q \) ist das Prinzip der Selbstorganisation.

\( Q \) ist die Quelle der Physik.

\section{Historischer Hintergrund: Reverse-Engineering der Physik}

Die IWT ist nicht aus dem Nichts entstanden. Sie ist das Ergebnis eines historischen Reverse-Engineering der gesamten Physik.

\begin{enumerate}
\item \textbf{Die Weber-Elektrodynamik (WED):} Eine feldlose, geschwindigkeitsabhängige Fernwirkungstheorie, die historisch älter ist als die Maxwell-Theorie. Sie erwies sich als fundamentaler als Maxwell – Maxwell ist der Grenzfall der WED.

\item \textbf{Die Weber-Gravitation (WG):} Die Übertragung der WED-Struktur auf die Gravitation. Die WG reproduziert alle Erfolge der ART (Periheldrehung, Lichtablenkung, Shapiro-Laufzeit) – jedoch ohne Raumzeitkrümmung.

\item \textbf{Die De-Broglie-Bohm-Theorie (DBT):} Die Erkenntnis, dass die mathematische Struktur der WED und WG identisch ist mit der Struktur der DBT. Alle drei sind nicht-lokal, instantan und deterministisch.

\item \textbf{Die Dynamische Schwere-Trägheits-Theorie (DSTT):} Eine eigenständige Theorie, die zeigt, dass das Äquivalenzprinzip nicht gilt. Die DSTT-Struktur ist in der WG und WED enthalten.

\item \textbf{Die Omega-Theorie:} Die Erkenntnis, dass nicht Krümmung, sondern Rotation das natürliche Prinzip der Gravitation ist.

\item \textbf{Die Maxwell-Gravitation:} Die Herleitung der Maxwell-Gleichungen der Gravitation aus der Wirbelstruktur der WG – und damit die natürliche Erklärung von Gravitationswellen mit Dispersion.

\item \textbf{Die IWT:} Die Rückführung aller dieser Theorien auf eine einzige Quelle – die Information und ihre Selbstorganisation \( Q \).
\end{enumerate}

Die gesamte Kette ist logisch und historisch konsistent. Die IWT ist der Abschluss dieser Kette.

\section{Überblick über die Arbeit}

Diese Arbeit ist in vier Teile gegliedert:

\begin{description}
\item[Teil I:] Die fundamentale Ebene – die Informations-Weber-Theorie (IWT) mit ihren Axiomen und der Herleitung von \( Q \).
\item[Teil II:] Die Emergenz der Physik aus \( Q \) – DBT, WG, WED, DSTT, Omega-Theorie, Maxwell-Gravitation.
\item[Teil III:] Die Bestätigung – die numerische Simulation der IWT und ihre Ergebnisse.
\item[Teil IV:] Die Folgen – die vereinheitlichte Physik, testbare Vorhersagen und die Schlussbetrachtung.
\end{description}

Anhänge enthalten die mathematischen Grundlagen, vollständige Herleitungen und einen Vergleich mit etablierten Theorien.